Single-member district systems face a foundational normative tension: should a party winning
$v\%$ of votes win approximately $v\%$ of seats (proportionality), or should the plurality
winner receive a seat-share bonus that magnifies its mandate (majoritarianism)? Neither standard
is constitutionally required under federal law post-\textit{Rucho}. But state constitutional
provisions---Pennsylvania's ``free and equal elections'' clause (Art.\ I, \S 5) and North
Carolina's ``free elections'' clause (Art.\ I, \S 10)---may impose a softer proportionality
floor that permits courts to invalidate maps deviating ``far beyond what any neutral criteria
would produce.''

This paper applies the Gallagher Least-Squares Index ($\text{LSq} = \sqrt{\tfrac{1}{2}
\sum_i(v_i - s_i)^2}$) to measure disproportionality for \textsc{bisect} algorithmic maps
and enacted 2022 congressional maps across NC ($k=14$), WI ($k=8$), TX ($k=38$), and
FL ($k=28$). \textsc{Bisect} maps achieve $\text{LSq} \in [0.04, 0.07]$---significantly
below enacted maps ($[0.10, 0.18]$) but not zero, because single-member districts
structurally preclude perfect proportionality: the majoritarian seat bonus ($1.3$--$1.5$x
for bisect vs.\ $2.1$x for enacted NC) is inherent in geographic representation with
geographic sorting of partisan voters.

The normative contribution is a principled answer to the defendant's most common response
to partisan gerrymandering claims---``proportionality is not required''---that (a) confirms
the claim is correct at the federal level post-\textit{Rucho}; (b) shows that state
constitutional text may impose a softer proportionality floor; and (c) demonstrates that
regardless of which normative standard governs, \textsc{bisect} maps satisfy both more
closely than enacted maps. Expert witnesses can frame the normative choice as a legal
question for the court while providing empirical evidence relevant under either standard.
